\section{Smooth ring maps}
\label{section-smooth}

\noindent
Let us motivate the definition of a smooth ring map by an example.
Suppose $R$ is a ring and $S = R[x, y]/(f)$ for some nonzero $f \in R[x, y]$.
In this case there is an exact sequence
$$
S \to
S\text{d}x \oplus S\text{d}y \to
\Omega_{S/R} \to 0
$$
where the first arrow maps $1$ to
$\frac{\partial f}{\partial x} \text{d}x +
\frac{\partial f}{\partial y} \text{d}y$ see
Section \ref{section-netherlander}.
We conclude that $\Omega_{S/R}$ is locally free of rank $1$ if
the partial derivatives of $f$ generate the unit ideal in $S$.
In this case $S$ is smooth of relative dimension $1$ over $R$.
But it can happen that $\Omega_{S/R}$ is locally free of rank $2$
namely if both partial derivatives of $f$ are zero. For example if
for a prime $p$ we have $p = 0$ in $R$ and $f = x^p + y^p$ then this
happens. Here $R \to S$ is a relative global complete intersection
of relative dimension $1$ which is not smooth.
Hence, in order to check that a ring map
is smooth it is not sufficient to check whether the module of differentials
is free. The correct condition is the following.

\begin{definition}
\label{definition-smooth}
A ring map $R \to S$ is {\it smooth} if it is of finite presentation
and the naive cotangent complex $\NL_{S/R}$ is quasi-isomorphic to a
finite projective $S$-module placed in degree $0$: this means
that $H_1(\NL_{S/R}) = 0$ and that $\Omega_{S/R}$ is a finite projective
$S$-module.
\end{definition}

\noindent
Let $R \to S$ be a ring map. By Lemma \ref{lemma-NL-homotopy} for any
presentation $\alpha : P \to S$ we have $H_1(\NL(\alpha)) = H_1(\NL_{S/R})$
and $H_0(\NL(\alpha)) = \Omega_{S/R}$. Thus, if $R \to S$ is smooth,
then for any surjection $\alpha : R[x_1, \ldots, x_n] \to S$ with kernel
$I$ the map
$$
I/I^2
\longrightarrow
\Omega_{R[x_1, \ldots, x_n]/R} \otimes_{R[x_1, \ldots, x_n]} S
$$
is an injective map whose cokernel is a finite projective $S$-module.
Thus the displayed arrow is a split injection. In other words
$\bigoplus_{i = 1}^n S \text{d}x_i \cong I/I^2 \oplus \Omega_{S/R}$
as $S$-modules and $I/I^2$ is a finite projective
$S$-module as well. Conversely, if $R \to S$ is of finite presentation
and we have a surjection $\alpha : R[x_1, \ldots, x_n] \to S$ with kernel
$I$ such that the displayed arrow is a split injection, then $R \to S$
is smooth.

\begin{lemma}
\label{lemma-localize-smooth}
Let $R \to S$ be a smooth ring map.
Any localization $S_g$ is smooth over $R$.
If $f \in R$ maps to an invertible element of $S$,
then $R_f \to S$ is smooth.
\end{lemma}

\begin{proof}
By Lemma \ref{lemma-localize-NL} the naive cotangent
complex for $S_g$ over $R$ is the base change of the naive cotangent
complex of $S$ over $R$. The assumption is that the naive cotangent
complex of $S/R$ is $\Omega_{S/R}$ and that this is a finite projective
$S$-module. Hence so is its base change. Thus $S_g$ is smooth over $R$.

\medskip\noindent
The second assertion follows in the same way from
Lemma \ref{lemma-NL-localize-bottom}.
\end{proof}

\begin{lemma}
\label{lemma-base-change-smooth}
\begin{slogan}
Smoothness is preserved under base change
\end{slogan}
Let $R \to S$ be a smooth ring map.
Let $R \to R'$ be any ring map.
Then the base change $R' \to S' = R' \otimes_R S$ is smooth.
\end{lemma}

\begin{proof}
Let $\alpha : R[x_1, \ldots, x_n] \to S$ be a presentation
with kernel $I$. Let $\alpha' : R'[x_1, \ldots, x_n] \to R' \otimes_R S$
be the induced presentation. Let $I' = \Ker(\alpha')$.
Since $0 \to I \to R[x_1, \ldots, x_n] \to S \to 0$
is exact, the sequence
$R' \otimes_R I \to R'[x_1, \ldots, x_n] \to R' \otimes_R S \to 0$
is exact. Thus $R' \otimes_R I \to I'$ is surjective.
By Definition \ref{definition-smooth} there is a short exact sequence
$$
0 \to I/I^2 \to
\Omega_{R[x_1, \ldots, x_n]/R} \otimes_{R[x_1, \ldots, x_n]} S \to
\Omega_{S/R} \to
0
$$
and the $S$-module $\Omega_{S/R}$ is finite projective.
In particular $I/I^2$ is a direct summand of
$\Omega_{R[x_1, \ldots, x_n]/R} \otimes_{R[x_1, \ldots, x_n]} S$.
Consider the commutative diagram
$$
\xymatrix{
R' \otimes_R (I/I^2) \ar[r] \ar[d] &
R' \otimes_R (\Omega_{R[x_1, \ldots, x_n]/R} \otimes_{R[x_1, \ldots, x_n]} S)
\ar[d] \\
I'/(I')^2 \ar[r] &
\Omega_{R'[x_1, \ldots, x_n]/R'}
\otimes_{R'[x_1, \ldots, x_n]} (R' \otimes_R S)
}
$$
Since the right vertical map is an isomorphism we see that
the left vertical map is injective and surjective by what was
said above. Thus we conclude that $\NL(\alpha')$ is quasi-isomorphic
to $\Omega_{S'/R'} \cong S' \otimes_S \Omega_{S/R}$ placed in degree $0$.
This module is finite projective since it is the base change
of a finite projective module.
\end{proof}

\begin{lemma}
\label{lemma-smooth-over-field}
Let $k$ be a field.
Let $S$ be a smooth $k$-algebra.
Then $S$ is a local complete intersection.
\end{lemma}

\begin{proof}
By Lemmas \ref{lemma-base-change-smooth} and
\ref{lemma-lci-field-change} it suffices to prove this when
$k$ is algebraically closed. Choose a presentation
$\alpha : k[x_1, \ldots, x_n] \to S$ with kernel $I$. Let $\mathfrak m$
be a maximal ideal of $S$, and let $\mathfrak m' \supset I$ be the
corresponding maximal ideal of $k[x_1, \ldots, x_n]$.
We will show that condition (5) of
Lemma \ref{lemma-lci}
holds (with $\mathfrak m$ instead of $\mathfrak q$).
We may write $\mathfrak m' = (x_1 - a_1, \ldots, x_n - a_n)$
for some $a_i \in k$, because $k$ is algebraically closed, see
Theorem \ref{theorem-nullstellensatz}.
By our assumption that $k \to S$ is smooth the $S$-module map
$\text{d} : I/I^2 \to \bigoplus_{i = 1}^n S \text{d}x_i$
is a split injection. Hence the corresponding map
$I/\mathfrak m' I \to \bigoplus \kappa(\mathfrak m') \text{d}x_i$
is injective. Say $\dim_{\kappa(\mathfrak m')}(I/\mathfrak m' I) = c$
and pick $f_1, \ldots, f_c \in I$ which map to a $\kappa(\mathfrak m')$-basis
of $I/\mathfrak m' I$. By
Nakayama's Lemma \ref{lemma-NAK}
we see that $f_1, \ldots, f_c$ generate $I_{\mathfrak m'}$ over
$k[x_1, \ldots, x_n]_{\mathfrak m'}$. Consider the commutative diagram
$$
\xymatrix{
I \ar[r] \ar[d] & I/I^2 \ar[rr] \ar[d] & &
I/\mathfrak m'I \ar[d] \\
\Omega_{k[x_1, \ldots, x_n]/k} \ar[r] &
\bigoplus S\text{d}x_i \ar[rr]^{\text{d}x_i \mapsto x_i - a_i} & &
\mathfrak m'/(\mathfrak m')^2
}
$$
(proof commutativity omitted). The middle vertical map is the one defining
the naive cotangent complex of $\alpha$. Note that the right lower
horizontal arrow induces an isomorphism
$\bigoplus \kappa(\mathfrak m') \text{d}x_i \to \mathfrak m'/(\mathfrak m')^2$.
Hence our generators $f_1, \ldots, f_c$ of $I_{\mathfrak m'}$ map to a
collection of elements in $k[x_1, \ldots, x_n]_{\mathfrak m'}$ whose
classes in $\mathfrak m'/(\mathfrak m')^2$ are linearly independent
over $\kappa(\mathfrak m')$. Therefore they form a regular sequence
in the ring $k[x_1, \ldots, x_n]_{\mathfrak m'}$ by
Lemma \ref{lemma-regular-ring-CM}.
This verifies condition (5) of
Lemma \ref{lemma-lci}
hence $S_g$ is a global complete intersection over $k$ for some
$g \in S$, $g \not \in \mathfrak m$. As this works for any maximal
ideal of $S$ we conclude that $S$ is a local complete intersection over $k$.
\end{proof}

\begin{definition}
\label{definition-standard-smooth}
Let $R$ be a ring. Given integers $n \geq c \geq 0$ and
$f_1, \ldots, f_c \in R[x_1, \ldots, x_n]$ we say
$$
R[x_1, \ldots, x_n]/(f_1, \ldots, f_c)
$$
is a {\it standard smooth algebra over $R$} if the polynomial
$$
g =
\det
\left(
\begin{matrix}
\partial f_1/\partial x_1 &
\partial f_2/\partial x_1 &
\ldots &
\partial f_c/\partial x_1 \\
\partial f_1/\partial x_2 &
\partial f_2/\partial x_2 &
\ldots &
\partial f_c/\partial x_2 \\
\ldots & \ldots & \ldots & \ldots \\
\partial f_1/\partial x_c &
\partial f_2/\partial x_c &
\ldots &
\partial f_c/\partial x_c
\end{matrix}
\right)
$$
maps to an invertible element in $R[x_1, \ldots, x_n]/(f_1, \ldots, f_c)$.
We say an $R$-algebra $S$ is {\it standard smooth} or that
the ring map $R \to S$ is {\it standard smooth} if there exist
$n \geq c \geq 0$ and $f_1, \ldots, f_c \in R[x_1, \ldots, x_n]$
such that $R[x_1, \ldots, x_n]/(f_1, \ldots, f_c)$ is a standard smooth
algebra over $R$ and $S$ is isomorphic to
$R[x_1, \ldots, x_n]/(f_1, \ldots, f_c)$ as an $R$-algebra.
\end{definition}

\begin{lemma}
\label{lemma-standard-smooth}
Let
$S = R[x_1, \ldots, x_n]/(f_1, \ldots, f_c) = R[x_1, \ldots, x_n]/I$
be a standard smooth algebra. Then
\begin{enumerate}
\item the ring map $R \to S$ is smooth,
\item the $S$-module $\Omega_{S/R}$ is free on
$\text{d}x_{c + 1}, \ldots, \text{d}x_n$,
\item the $S$-module $I/I^2$ is free on the classes of $f_1, \ldots, f_c$,
\item for any $g \in S$ the ring map $R \to S_g$ is standard smooth,
\item for any ring map $R \to R'$ the base change
$R' \to R'\otimes_R S$ is standard smooth,
\item if $f \in R$ maps to an invertible element in $S$, then
$R_f \to S$ is standard smooth, and
\item the ring $S$ is a relative global complete intersection over $R$.
\end{enumerate}
\end{lemma}

\begin{proof}
Consider the naive cotangent complex of the given presentation
$$
(f_1, \ldots, f_c)/(f_1, \ldots, f_c)^2
\longrightarrow
\bigoplus\nolimits_{i = 1}^n S \text{d}x_i
$$
Let us compose this map with the projection onto the first $c$ direct summands
of the direct sum. According to the definition of a standard smooth
algebra the classes $f_i \bmod (f_1, \ldots, f_c)^2$ map to a basis of
$\bigoplus_{i = 1}^c S\text{d}x_i$. We conclude that
$(f_1, \ldots, f_c)/(f_1, \ldots, f_c)^2$ is free of rank $c$ with
a basis given by the elements $f_i \bmod (f_1, \ldots, f_c)^2$, and
that the homology in degree $0$, i.e., $\Omega_{S/R}$,
of the naive cotangent complex is a free $S$-module with basis the images of
$\text{d}x_{c + j}$, $j = 1, \ldots, n - c$.
In particular, this proves $R \to S$ is smooth.

\medskip\noindent
The proofs of (4) and (6) are omitted. But see the example below and
the proof of
Lemma \ref{lemma-base-change-relative-global-complete-intersection}.

\medskip\noindent
Let $\varphi : R \to R'$ be any ring map.
Denote $S' = R'[x_1, \ldots, x_n]/(f_1^\varphi, \ldots, f_c^\varphi)$
where $f^\varphi$ is the polynomial obtained from $f \in R[x_1, \ldots, x_n]$
by applying $\varphi$ to all the coefficients. Then $S' \cong R' \otimes_R S$.
Moreover, the determinant of Definition \ref{definition-standard-smooth}
for $S'/R'$ is equal to $g^\varphi$. Its image in $S'$ is therefore
the image of $g$ via $R[x_1, \ldots, x_n] \to S \to S'$
and hence invertible. This proves (5).

\medskip\noindent
To prove (7) it suffices to show that
$S \otimes_R \kappa(\mathfrak p)$ has dimension $n - c$
for every prime $\mathfrak p \subset R$.
By (5) it suffices to prove that any standard smooth
algebra $k[x_1, \ldots, x_n]/(f_1, \ldots, f_c)$
over a field $k$ has dimension $n - c$. We already
know that $k[x_1, \ldots, x_n]/(f_1, \ldots, f_c)$ is a local
complete intersection by Lemma \ref{lemma-smooth-over-field}.
Hence, since $I/I^2$ is free of rank $c$ we see that
$k[x_1, \ldots, x_n]/(f_1, \ldots, f_c)$ has dimension
$n - c$, by Lemma \ref{lemma-lci} for example.
\end{proof}

\begin{example}
\label{example-make-standard-smooth}
Let $R$ be a ring.
Let $f_1, \ldots, f_c \in R[x_1, \ldots, x_n]$.
Let
$$
h =
\det
\left(
\begin{matrix}
\partial f_1/\partial x_1 &
\partial f_2/\partial x_1 &
\ldots &
\partial f_c/\partial x_1 \\
\partial f_1/\partial x_2 &
\partial f_2/\partial x_2 &
\ldots &
\partial f_c/\partial x_2 \\
\ldots & \ldots & \ldots & \ldots \\
\partial f_1/\partial x_c &
\partial f_2/\partial x_c &
\ldots &
\partial f_c/\partial x_c
\end{matrix}
\right).
$$
Set $S = R[x_1, \ldots, x_{n + 1}]/(f_1, \ldots, f_c, x_{n + 1}h - 1)$.
This is an example of a standard smooth algebra, except that the
presentation is wrong and the variables should be in the
following order:
$x_1, \ldots, x_c, x_{n + 1}, x_{c + 1}, \ldots, x_n$.
\end{example}

\begin{lemma}
\label{lemma-compose-standard-smooth}
A composition of standard smooth ring maps is standard smooth.
\end{lemma}

\begin{proof}
Suppose that $R \to S$ and $S \to S'$ are standard smooth. We choose
presentations
$S =  R[x_1, \ldots, x_n]/(f_1, \ldots, f_c)$
and
$S' = S[y_1, \ldots, y_m]/(g_1, \ldots, g_d)$.
Choose elements $g_j' \in R[x_1, \ldots, x_n, y_1, \ldots, y_m]$ mapping
to the $g_j$. In this way we see
$S' = R[x_1, \ldots, x_n, y_1, \ldots, y_m]/
(f_1, \ldots, f_c, g'_1, \ldots, g'_d)$.
To show that $S'$ is standard smooth it suffices to verify
that the determinant
$$
\det
\left(
\begin{matrix}
\partial f_1/\partial x_1 &
\ldots &
\partial f_c/\partial x_1 &
\partial g_1/\partial x_1 &
\ldots &
\partial g_d/\partial x_1 \\
\ldots &
\ldots &
\ldots &
\ldots &
\ldots &
\ldots  \\
\partial f_1/\partial x_c &
\ldots &
\partial f_c/\partial x_c &
\partial g_1/\partial x_c &
\ldots &
\partial g_d/\partial x_c \\
0 &
\ldots &
0 &
\partial g_1/\partial y_1 &
\ldots &
\partial g_d/\partial y_1 \\
\ldots &
\ldots &
\ldots &
\ldots &
\ldots &
\ldots  \\
0 &
\ldots &
0 &
\partial g_1/\partial y_d &
\ldots &
\partial g_d/\partial y_d
\end{matrix}
\right)
$$
is invertible in $S'$. This is clear since it is the product
of the two determinants which were assumed to be invertible
by hypothesis.
\end{proof}

\begin{lemma}
\label{lemma-smooth-syntomic}
Let $R \to S$ be a smooth ring map.
There exists an open covering of $\Spec(S)$ by
standard opens $D(g)$ such that each $S_g$ is standard smooth
over $R$. In particular $R \to S$ is syntomic.
\end{lemma}

\begin{proof}
Choose a presentation $\alpha : R[x_1, \ldots, x_n] \to S$
with kernel $I = (f_1, \ldots, f_m)$. For every subset
$E \subset \{1, \ldots, m\}$ consider the open
subset $U_E$ where the classes $f_e, e\in E$ freely generate
the finite projective $S$-module $I/I^2$, see Lemma \ref{lemma-cokernel-flat}.
We may cover $\Spec(S)$ by standard opens $D(g)$ each
completely contained in one of the opens $U_E$. For such a $g$
we look at the presentation
$$
\beta : R[x_1, \ldots, x_n, x_{n + 1}] \longrightarrow S_g
$$
mapping $x_{n + 1}$ to $1/g$. Setting $J = \Ker(\beta)$ we
use Lemma \ref{lemma-principal-localization-NL} to see that
$J/J^2 \cong (I/I^2)_g \oplus S_g$ is free.
We may and do replace $S$ by $S_g$. Then using
Lemma \ref{lemma-huber} we may assume we have a presentation
$\alpha : R[x_1, \ldots, x_n] \to S$ with kernel $I = (f_1, \ldots, f_c)$
such that $I/I^2$ is free on the classes of $f_1, \ldots, f_c$.

\medskip\noindent
Using the presentation $\alpha$ obtained at the end of the previous
paragraph, we more or less repeat this argument with
the basis elements $\text{d}x_1, \ldots, \text{d}x_n$
of $\Omega_{R[x_1, \ldots, x_n]/R}$.
Namely, for any subset $E \subset \{1, \ldots, n\}$ of cardinality $c$
we may consider the open subset $U_E$ of $\Spec(S)$ where
the differential of $\NL(\alpha)$ composed with the projection
$$
S^{\oplus c} \cong I/I^2
\longrightarrow
\Omega_{R[x_1, \ldots, x_n]/R} \otimes_{R[x_1, \ldots, x_n]} S
\longrightarrow
\bigoplus\nolimits_{i \in E} S\text{d}x_i
$$
is an isomorphism. Again we may find a covering of $\Spec(S)$
by (finitely many) standard opens $D(g)$ such that each $D(g)$
is completely contained in one of the opens $U_E$.
By renumbering, we may assume $E = \{1, \ldots, c\}$.
For a $g$ with $D(g) \subset U_E$ we look at the presentation
$$
\beta : R[x_1, \ldots, x_n, x_{n + 1}] \to S_g
$$
mapping $x_{n + 1}$ to $1/g$. Setting $J = \Ker(\beta)$
we conclude from Lemma \ref{lemma-principal-localization-NL}
that $J = (f_1, \ldots, f_c, fx_{n + 1} - 1)$ where $\alpha(f) = g$
and that the composition
$$
J/J^2 \longrightarrow
\Omega_{R[x_1, \ldots, x_{n + 1}]/R} \otimes_{R[x_1, \ldots, x_{n + 1}]} S_g
\longrightarrow
\bigoplus\nolimits_{i = 1}^c S_g\text{d}x_i \oplus S_g \text{d}x_{n + 1}
$$
is an isomorphism. Reordering the coordinates as
$x_1, \ldots, x_c, x_{n + 1}, x_{c + 1}, \ldots, x_n$
we conclude that $S_g$ is standard smooth over $R$ as desired.

\medskip\noindent
This finishes the proof as standard smooth algebras are syntomic
(Lemmas \ref{lemma-standard-smooth} and
\ref{lemma-relative-global-complete-intersection})
and being syntomic over $R$ is local on $S$
(Lemma \ref{lemma-local-syntomic}).
\end{proof}

\begin{definition}
\label{definition-smooth-at-prime}
Let $R \to S$ be a ring map.
Let $\mathfrak q$ be a prime of $S$.
We say $R \to S$ is {\it smooth at $\mathfrak q$} if there
exists a $g \in S$, $g \not \in \mathfrak q$ such
that $R \to S_g$ is smooth.
\end{definition}

\noindent
For ring maps of finite presentation we can characterize this as follows.

\begin{lemma}
\label{lemma-smooth-at-point}
Let $R \to S$ be of finite presentation. Let $\mathfrak q$ be a
prime of $S$. The following are equivalent
\begin{enumerate}
\item $R \to S$ is smooth at $\mathfrak q$,
\item $H_1(L_{S/R})_\mathfrak q = 0$ and
$\Omega_{S/R, \mathfrak q}$ is a finite free $S_\mathfrak q$-module,
\item $H_1(L_{S/R})_\mathfrak q = 0$ and
$\Omega_{S/R, \mathfrak q}$ is a projective $S_\mathfrak q$-module, and
\item $H_1(L_{S/R})_\mathfrak q = 0$ and
$\Omega_{S/R, \mathfrak q}$ is a flat $S_\mathfrak q$-module.
\end{enumerate}
\end{lemma}

\begin{proof}
We will use without further mention that formation of the
naive cotangent complex commutes with localization, see
Section \ref{section-netherlander}, especially
Lemma \ref{lemma-localize-NL}.
Note that $\Omega_{S/R}$ is a finitely presented $S$-module, see
Lemma \ref{lemma-differentials-finitely-presented}. Hence
(2), (3), and (4) are equivalent by Lemma \ref{lemma-finite-projective}.
It is clear that (1) implies the equivalent conditions (2), (3), and (4).
Assume (2) holds. Writing $S_\mathfrak q$ as
the colimit of principal localizations we see from
Lemma \ref{lemma-colimit-category-fp-modules}
that we can find a $g \in S$, $g \not \in \mathfrak q$ such that
$(\Omega_{S/R})_g$ is finite free. Choose a presentation
$\alpha : R[x_1, \ldots, x_n] \to S$ with kernel $I$. We may work
with $\NL(\alpha)$ instead of $\NL_{S/R}$, see
Lemma \ref{lemma-NL-homotopy}. The surjection
$$
\Omega_{R[x_1, \ldots, x_n]/R} \otimes_{R[x_1, \ldots, x_n]} S
\to \Omega_{S/R} \to 0
$$
has a right inverse after inverting $g$ because $(\Omega_{S/R})_g$ is
projective. Hence the image of
$\text{d} : (I/I^2)_g \to
\Omega_{R[x_1, \ldots, x_n]/R} \otimes_{R[x_1, \ldots, x_n]} S_g$
is a direct summand and this map has a right inverse too.
We conclude that $H_1(L_{S/R})_g$ is a quotient of $(I/I^2)_g$.
In particular $H_1(L_{S/R})_g$ is a finite $S_g$-module.
Thus the vanishing of $H_1(L_{S/R})_{\mathfrak q}$
implies the vanishing of $H_1(L_{S/R})_{gg'}$ for some $g' \in S$,
$g' \not \in \mathfrak q$. Then $R \to S_{gg'}$ is smooth by
definition.
\end{proof}

\begin{lemma}
\label{lemma-locally-smooth}
\begin{slogan}
A ring map is smooth if and only if it is smooth at all primes of the target
\end{slogan}
Let $R \to S$ be a ring map.
Then $R \to S$ is smooth if and only if $R \to S$ is smooth
at every prime $\mathfrak q$ of $S$.
\end{lemma}

\begin{proof}
The direct implication is trivial. Suppose that $R \to S$ is smooth
at every prime $\mathfrak q$ of $S$. Since $\Spec(S)$ is
quasi-compact, see Lemma \ref{lemma-quasi-compact},
there exists a finite covering
$\Spec(S) = \bigcup D(g_i)$ such that each $S_{g_i}$ is
smooth. By Lemma \ref{lemma-cover-upstairs} this implies that
$S$ is of finite presentation over $R$. According to
Lemma \ref{lemma-localize-NL} we see that
$\NL_{S/R} \otimes_S S_{g_i}$ is quasi-isomorphic to a finite projective
$S_{g_i}$-module placed in degree $0$. By Lemma \ref{lemma-finite-projective}
this implies that $\NL_{S/R}$ is quasi-isomorphic to a finite
projective $S$-module placed in degree $0$.
\end{proof}

\begin{lemma}
\label{lemma-compose-smooth}
A composition of smooth ring maps is smooth.
\end{lemma}

\begin{proof}
You can prove this in many different ways. One way is to use
the snake lemma (Lemma \ref{lemma-snake}), the Jacobi-Zariski sequence
(Lemma \ref{lemma-exact-sequence-NL}), combined with the
characterization of projective modules as being
direct summands of free modules (Lemma \ref{lemma-characterize-projective}).
Another proof can be obtained by combining
Lemmas \ref{lemma-smooth-syntomic}, \ref{lemma-compose-standard-smooth}
and \ref{lemma-locally-smooth}.
\end{proof}

\begin{lemma}
\label{lemma-product-smooth}
Let $R$ be a ring. Let $S = S' \times S''$ be a product of $R$-algebras.
Then $S$ is smooth over $R$ if and only if both $S'$ and $S''$ are
smooth over $R$.
\end{lemma}

\begin{proof}
Omitted. Hints: By Lemma \ref{lemma-locally-smooth} we can check smoothness
one prime at a time. Since $\Spec(S)$ is the disjoint union of
$\Spec(S')$ and $\Spec(S'')$ by Lemma \ref{lemma-spec-product}
we find that smoothness of $R \to S$ at $\mathfrak q$ corresponds
to either smoothness of $R \to S'$ at the corresponding prime or
smoothness of $R \to S''$ at the corresponding prime.
\end{proof}

\begin{lemma}
\label{lemma-relative-global-complete-intersection-smooth}
Let $R$ be a ring. Let $S = R[x_1, \ldots, x_n]/(f_1, \ldots, f_c)$
be a relative global complete intersection.
Let $\mathfrak q \subset S$ be a prime. Then $R \to S$
is smooth at $\mathfrak q$ if and only if there exists a
subset $I \subset \{1, \ldots, n\}$ of cardinality $c$
such that the polynomial
$$
g_I = \det (\partial f_j/\partial x_i)_{j = 1, \ldots, c, \ i \in I}.
$$
does not map to an element of $\mathfrak q$.
\end{lemma}

\begin{proof}
By Lemma \ref{lemma-relative-global-complete-intersection-conormal}
we see that the naive cotangent complex
associated to the given presentation of $S$ is the complex
$$
\bigoplus\nolimits_{j = 1}^c S \cdot f_j
\longrightarrow
\bigoplus\nolimits_{i = 1}^n S \cdot \text{d}x_i, \quad
f_j \longmapsto \sum \frac{\partial f_j}{\partial x_i} \text{d}x_i.
$$
The maximal minors of the matrix giving the map are exactly
the polynomials $g_I$.

\medskip\noindent
Assume $g_I$ maps to $g \in S$, with $g \not \in \mathfrak q$.
Then the algebra $S_g$ is smooth over $R$. Namely, its naive
cotangent complex is quasi-isomorphic to the complex above
localized at $g$, see Lemma \ref{lemma-localize-NL}. And by
construction it is quasi-isomorphic to a free rank $n - c$
module in degree $0$.

\medskip\noindent
Conversely, suppose that all $g_I$ end up in $\mathfrak q$.
In this case the complex above tensored with $\kappa(\mathfrak q)$
does not have maximal rank, and hence there is no localization
by an element $g \in S$, $g \not \in \mathfrak q$
where this map becomes a split injection. By Lemma \ref{lemma-localize-NL}
again there is no such localization which is smooth over $R$.
\end{proof}

\begin{lemma}
\label{lemma-flat-fibre-smooth}
Let $R \to S$ be a ring map.
Let $\mathfrak q \subset S$ be a prime lying over the
prime $\mathfrak p$ of $R$. Assume
\begin{enumerate}
\item there exists a $g \in S$, $g \not\in \mathfrak q$
such that $R \to S_g$ is of finite presentation,
\item the local ring homomorphism
$R_{\mathfrak p} \to S_{\mathfrak q}$ is flat,
\item the fibre $S \otimes_R \kappa(\mathfrak p)$ is smooth
over $\kappa(\mathfrak p)$ at the prime corresponding
to $\mathfrak q$.
\end{enumerate}
Then $R \to S$ is smooth at $\mathfrak q$.
\end{lemma}

\begin{proof}
By Lemmas \ref{lemma-syntomic} and \ref{lemma-smooth-over-field}
we see that there exists a $g \in S$ such that $S_g$ is a
relative global complete intersection. Replacing $S$ by $S_g$ we may assume
$S = R[x_1, \ldots, x_n]/(f_1, \ldots, f_c)$ is a relative
global complete intersection.
For any subset $I \subset \{1, \ldots, n\}$ of cardinality
$c$ consider the polynomial
$g_I = \det (\partial f_j/\partial x_i)_{j = 1, \ldots, c, i \in I}$
of Lemma \ref{lemma-relative-global-complete-intersection-smooth}.
Note that the image $\overline{g}_I$ of $g_I$ in the polynomial ring
$\kappa(\mathfrak p)[x_1, \ldots, x_n]$ is the determinant
of the partial derivatives of the images $\overline{f}_j$ of the $f_j$
in the ring $\kappa(\mathfrak p)[x_1, \ldots, x_n]$. Thus the lemma follows
by applying Lemma \ref{lemma-relative-global-complete-intersection-smooth}
both to $R \to S$ and to
$\kappa(\mathfrak p) \to S \otimes_R \kappa(\mathfrak p)$.
\end{proof}

\noindent
Note that the sets $U, V$ in the following lemma
are open by definition.

\begin{lemma}
\label{lemma-flat-base-change-locus-smooth}
Let $R \to S$ be a ring map of finite presentation.
Let $R \to R'$ be a flat ring map.
Denote $S' = R' \otimes_R S$ the base change.
Let $U \subset \Spec(S)$ be the set of primes at
which $R \to S$ is smooth.
Let $V \subset \Spec(S')$ the set of primes at
which $R' \to S'$ is smooth.
Then $V$ is the inverse image of $U$ under the
map $f : \Spec(S') \to \Spec(S)$.
\end{lemma}

\begin{proof}
By Lemma \ref{lemma-change-base-NL} we see that
$\NL_{S/R} \otimes_S S'$ is homotopy equivalent to $\NL_{S'/R'}$.
This already implies that $f^{-1}(U) \subset V$.

\medskip\noindent
Let $\mathfrak q' \subset S'$ be a prime lying over
$\mathfrak q \subset S$. Assume $\mathfrak q' \in V$.
We have to show that $\mathfrak q \in U$.
Since $S \to S'$ is flat, we see that $S_{\mathfrak q} \to S'_{\mathfrak q'}$
is faithfully flat (Lemma \ref{lemma-local-flat-ff}). Thus the vanishing of
$H_1(L_{S'/R'})_{\mathfrak q'}$ implies the
vanishing of $H_1(L_{S/R})_{\mathfrak q}$.
By Lemma \ref{lemma-finite-projective-descends}
applied to the $S_{\mathfrak q}$-module $(\Omega_{S/R})_{\mathfrak q}$
and the map $S_{\mathfrak q} \to S'_{\mathfrak q'}$ we see that
$(\Omega_{S/R})_{\mathfrak q}$ is projective. Hence
$R \to S$ is smooth at $\mathfrak q$ by
Lemma \ref{lemma-smooth-at-point}.
\end{proof}

\begin{lemma}
\label{lemma-smooth-field-change-local}
Let $K/k$ be a field extension.
Let $S$ be a finite type algebra over $k$.
Let $\mathfrak q_K$ be a prime of $S_K = K \otimes_k S$
and let $\mathfrak q$ be the corresponding prime of $S$.
Then $S$ is smooth over $k$ at $\mathfrak q$ if and only if
$S_K$ is smooth at $\mathfrak q_K$ over $K$.
\end{lemma}

\begin{proof}
This is a special case of Lemma \ref{lemma-flat-base-change-locus-smooth}.
\end{proof}

\begin{lemma}
\label{lemma-lift-smooth}
Let $R$ be a ring and let $I \subset R$ be an ideal.
Let $R/I \to \overline{S}$ be a smooth ring map.
Then there exists elements $\overline{g}_i \in \overline{S}$
which generate the unit ideal of $\overline{S}$
such that each $\overline{S}_{g_i} \cong S_i/IS_i$
for some (standard) smooth ring $S_i$ over $R$.
\end{lemma}

\begin{proof}
By Lemma \ref{lemma-smooth-syntomic} we find a collection of elements
$\overline{g}_i \in \overline{S}$
which generate the unit ideal of $\overline{S}$
such that each $\overline{S}_{g_i}$ is standard smooth over $R/I$.
Hence we may assume that $\overline{S}$ is standard smooth
over $R/I$. Write
$\overline{S} =
(R/I)[x_1, \ldots, x_n]/(\overline{f}_1, \ldots, \overline{f}_c)$
as in Definition \ref{definition-standard-smooth}.
Choose $f_1, \ldots, f_c \in R[x_1, \ldots, x_n]$
lifting $\overline{f}_1, \ldots, \overline{f}_c$. Set
$S = R[x_1, \ldots, x_n, x_{n + 1}]/(f_1, \ldots, f_c, x_{n + 1}\Delta - 1)$
where $\Delta = \det(\frac{\partial f_j}{\partial x_i})_{i, j = 1, \ldots, c}$
as in Example \ref{example-make-standard-smooth}.
This proves the lemma.
\end{proof}








